% Figure: two cantilevers of uniform strength compared with a prismatic
% one. A cantilever under a tip load carries a moment M(x)=Px growing
% linearly from the tip. Holding the extreme-fibre stress constant
% requires the section modulus to grow with x: a constant-width beam
% needs its depth to grow as sqrt(x) (parabolic profile), a
% constant-depth beam needs its width to grow linearly (triangular plan).
% The prismatic beam wastes material near the tip.
\begin{figure}[htbp]
\centering
\begin{tikzpicture}[scale=1.0]
% --- left panel: constant width, depth grows as sqrt(x) (parabolic elevation)
\begin{scope}[shift={(0,0)}]
% wall
\draw[enggray, very thick] (3.2,-1.1) -- (3.2,1.1);
\foreach \y in {-0.9,-0.5,-0.1,0.3,0.7} \draw[enggray] (3.2,\y) -- (3.5,\y+0.2);
% parabolic depth profile: depth ~ sqrt(distance from tip); tip at x=0
\draw[engblue, very thick]
  (0,0) .. controls (1.0,0.30) and (2.1,0.52) .. (3.2,0.62)
  -- (3.2,-0.62)
  .. controls (2.1,-0.52) and (1.0,-0.30) .. (0,0) -- cycle;
% tip load
\draw[engred, very thick, ->] (0,0.95) -- (0,0.05);
\node[engred, anchor=south, font=\scriptsize] at (0,1.0) {$P$};
\node[engblack, font=\scriptsize, anchor=north] at (1.6,-0.8) {depth $\propto\sqrt{x}$};
\node[engblack, font=\footnotesize, anchor=south] at (1.6,0.72) {constant width};
\end{scope}
% --- right panel: constant depth, width grows linearly (triangular plan view)
\begin{scope}[shift={(4.6,0)}]
\draw[enggray, very thick] (3.2,-1.1) -- (3.2,1.1);
\foreach \y in {-0.9,-0.5,-0.1,0.3,0.7} \draw[enggray] (3.2,\y) -- (3.5,\y+0.2);
% triangular plan: width grows linearly from tip
\draw[engamber, very thick]
  (0,0) -- (3.2,0.7) -- (3.2,-0.7) -- cycle;
\draw[engred, very thick, ->] (0,0.95) -- (0,0.05);
\node[engred, anchor=south, font=\scriptsize] at (0,1.0) {$P$};
\node[engblack, font=\scriptsize, anchor=north] at (1.6,-0.8) {width $\propto x$};
\node[engblack, font=\footnotesize, anchor=south] at (1.6,0.55) {constant depth};
\end{scope}
\end{tikzpicture}
\caption[Cantilevers of uniform strength]{Two cantilevers of uniform
strength under a tip load, whose bending moment $M(x)=Px$ grows linearly
from the tip. Keeping the extreme-fibre stress $\sigma=M/S$ constant along
the span requires the section modulus to grow in step with the moment.
Left, a beam of constant width meets this by letting its depth grow as
$\sqrt{x}$, a parabolic elevation. Right, a beam of constant depth meets
it by letting its width grow linearly, a triangular plan, which is the
leaf spring's diamond stack seen from above. Either profile removes the
material a prismatic beam wastes near the lightly stressed tip; the
uniform-strength shape is the lightest section that keeps every fibre at
the same working stress, at the cost of a tip that carries the deflection
because it is the most flexible part of the member.}
\label{fig:vol03ch08:uniform-strength}
\end{figure}
